\chapter*{How to read this catalogue}
\addcontentsline{toc}{chapter}{How to read this catalogue}

This volume contains the proofs of the results stated in the paper \emph{Exact Recovery of a Linear Differential Operator from a Single Solution} (complete edition, Phases~0--10). It is a catalogue: every result has one card, and every card has the same shape.

\section*{The shape of a card}

\begin{dimcard}{A0}{ID-T0}{Name of the result}{\proved}{}
The full statement, with every hypothesis.
\end{dimcard}
\begin{proofidea}
One paragraph that says why the result is true, before any computation.
\end{proofidea}
\begin{proof}
\step{1}{First claim}
One claim per step, with its computation displayed on its own line.
\step{2}{Second claim}
\dots
\whyworks{one sentence that names the mechanism of the proof.}
\end{proof}
\uses{the results this proof relies on, by catalogue number.}

\section*{The four parts}

\textbf{Part~I} contains the 31 results of Phases~0--5 on which everything else rests (Dim.~A1--A31). They were selected from the project's dependency graph: a result is in Part~I if many other results depend on it, or if it is a headline theorem of the paper. Their proofs are given in full, step by step.

\textbf{Part~II} contains the load-bearing results of Phases~6--10 (Dim.~A32--A53): systems and gauge, variable coefficients, convergence and noise, the analytic encoder, the universal dictionary, and the Phase~10 corrections to results of Phases~0--5. They were selected by the same criterion: use in the extended algorithm (Chapter~\ref{PAP-ch:extalgo} of the paper), dependence in the ledger, and headline status.

\textbf{Part~III} contains all the other results of Phases~0--5, grouped by phase and ordered by identifier (Dim.~B1--B72). Short and medium proofs are given in full. Proofs that take many pages in the research files (the Burgers asymptotics, the scattering identification, the Hirota computations) are given as a condensed card, with every step and key formula and a pointer to the research file.

\textbf{Part~IV} contains the other results of Phases~6--10 (Dim.~B73--B155 and list cards). Eighty-three results have a card with a proof (Dim.~B73--B155); the others are \emph{list cards}: the statement, its trust label and the research file that holds the proof. All proofs of Part~II and Part~IV were rewritten from the research packages for this volume and were not re-audited; the trust label refers to the result as recorded in the ledger. This is a deliberate limit of this edition: for Phases~6--10 the research packages are about $150{,}000$ words, and only the proofs that the extended algorithm or a corrected statement depends on were rewritten here.

\textbf{Which edition to read.} The companion volumes for Phases~0--5 alone are \emph{paper\_fasi05} and \emph{dimostrazioni\_fasi05}. In case of a difference, this volume prevails: Phase~10 refuted or corrected several statements of Phases~0--5 (Chapter~\ref{ch:cIII6}).

\section*{Trust labels}

\begin{center}\small
\begin{tabular}{@{}L{4.4cm}L{10.0cm}@{}}
\toprule
Label & Meaning\\\midrule
\proved & complete proof, passed an independent review\\
\modulo{X} & complete except for a cited published theorem X, whose source was checked\\
\pw & complete proof, not independently reviewed\\
\cf & a published theorem, used as a black box\\
\condl & depends on a lemma that is not proved\\
\vex & checked on examples only\\
\falsest & a counterexample is proved (the card proves the counterexample)\\
\openq & unsolved\\\bottomrule
\end{tabular}
\end{center}

The reviewers were independent AI agents, not human referees. Identifiers such as C3-T1 are those of the project ledger \texttt{claims/LEDGER.md}; the research file that holds each original proof is named in the index at the end.

\section*{The backbone of Part~I (Phases~0--5)}

\begin{center}
\begin{tikzpicture}[x=1cm,y=1cm,
  n/.style={draw=cA,fill=cA!8,rounded corners=2pt,font=\scriptsize\sffamily,align=center,minimum height=8mm,text width=2.0cm},
  m/.style={draw=cB,fill=cB!8,rounded corners=2pt,font=\scriptsize\sffamily,align=center,minimum height=8mm,text width=2.0cm},
  a/.style={-{Stealth[length=1.6mm]},cG,thick}]
\node[n] (L2) at (0,0) {A1 gcd split\\C3-L2};
\node[n] (L3) at (0,-1.4) {A2 EP test\\C3-L3};
\node[n] (T1) at (3,-0.7) {A3 criterion\\C3-T1};
\node[n] (T6) at (6,0) {A4 $G_u$\\C3-T6};
\node[n] (P4) at (6,-1.4) {A5 spectral\\P4-T2};
\node[n] (T3) at (3,-2.8) {A6 thresholds\\C3-T2/T3};
\node[n] (P1) at (6,-2.8) {A7 no certificate\\C3-P1};
\node[n] (T4) at (9,-1.4) {A8 decidable\\P4-T4};
\node[m] (T5) at (9,-2.8) {A9 black box\\P4-T5};
\node[n] (CL3) at (0,-4.4) {A11 quotient\\C12-L3};
\node[n] (L1p) at (0,-5.8) {A12 twisted\\Poincar\'e};
\node[n] (T1p) at (3,-5.1) {A13 factors\\C12-T1$'$};
\node[n] (P7) at (6,-5.1) {A17 $\gamma_N$\\C12-P7};
\node[n] (P1T2) at (9,-5.1) {A19 $\Eint$\\P1-T2};
\node[n] (S4) at (9,-6.5) {A29 dictionary\\S4ID-T1};
\node[n] (J) at (6,-6.5) {A25 jet\\S4T-T1};
\draw[a] (L2)--(T1); \draw[a] (L3)--(T1); \draw[a] (T1)--(T6); \draw[a] (T1)--(P4);
\draw[a] (L3)--(T3); \draw[a] (T3)--(T4); \draw[a] (P4)--(T4); \draw[a] (T1)--(P1); \draw[a] (P1)--(T5);
\draw[a] (CL3)--(T1p); \draw[a] (L1p)--(T1p); \draw[a] (T1p)--(P7); \draw[a] (T1p)--(P1T2);
\draw[a] (P1T2)--(J); \draw[a] (L3)--(CL3);
\draw[a] (J)--(S4);
\end{tikzpicture}
\end{center}
{\small The arrows mean ``is used in the proof of''. Blue: positive results; red: a negative (impossibility) result. The arrows come from the \texttt{depends\_on} edges of the project graph.}
